%% AIML Engineering Lab - Project Report (v5)
%% Placeholders are replaced by generate_content.py
\documentclass[11pt,a4paper]{article}

\usepackage[a4paper, margin=2cm, top=1.5cm, bottom=1.5cm]{geometry}
\usepackage[T1]{fontenc}
\usepackage[utf8]{inputenc}
\usepackage{lmodern}
\usepackage{microtype}
\usepackage{setspace}
\usepackage{graphicx}
\usepackage{xcolor}
\usepackage{colortbl}
\usepackage{amsmath}
\usepackage{booktabs}
\usepackage{tabularx}
\usepackage{array}
\usepackage{enumitem}
\usepackage{fancyhdr}
\usepackage{titlesec}
\usepackage{hyperref}
\usepackage{float}
\usepackage{caption}
\usepackage{tcolorbox}
\usepackage{needspace}
\tcbuselibrary{skins,breakable}

%% ── Colour palette ─────────────────────────────────────────────────────────
\definecolor{primary}{HTML}{1A1A2E}
\definecolor{accent}{HTML}{0F3460}
\definecolor{titleblue}{HTML}{2E86C1}
\definecolor{headingblue}{HTML}{2196F3}
\definecolor{subheadblue}{HTML}{42A5F5}
\definecolor{codemuted}{HTML}{6B7280}
\definecolor{metricgreen}{HTML}{16A34A}
\definecolor{tableheader}{HTML}{D0E8F7}
\definecolor{tablerow}{HTML}{EBF5FB}
\definecolor{lightblue}{HTML}{F0F7FC}

\hypersetup{
  colorlinks=true,
  linkcolor=accent,
  urlcolor=accent,
  citecolor=accent,
  pdftitle={Hyperparameter Optimization},
  pdfauthor={AIML Engineering Lab},
}

\setstretch{1.08}
\setlength{\parskip}{3pt}
\setlength{\parindent}{0pt}

\titleformat{\section}
  {\large\bfseries\color{headingblue}}
  {\thesection}{1em}{}[\color{tableheader}\titlerule]
\titleformat{\subsection}
  {\normalsize\bfseries\color{subheadblue}}
  {\thesubsection}{0.8em}{}
\titlespacing{\section}{0pt}{8pt}{3pt}
\titlespacing{\subsection}{0pt}{6pt}{2pt}

\pagestyle{fancy}
\fancyhf{}
\fancyhead[L]{\small\color{codemuted}Project 010}
\fancyhead[R]{\small\color{codemuted}Hyperparameter Optimization}
\fancyfoot[C]{\small\color{codemuted}\thepage}
\renewcommand{\headrulewidth}{0.2pt}

\newtcolorbox{resultbox}{
  colback=lightblue, colframe=tableheader,
  boxrule=0.4pt, arc=4pt, left=10pt, right=10pt, top=3pt, bottom=3pt
}

\begin{document}

%% ── PAGE 1: Title + Overview ───────────────────────────────────────────────
\thispagestyle{empty}

{\colorbox{primary}{\makebox[\dimexpr\linewidth-2\fboxsep][l]{\color{white}\bfseries\normalsize\hspace{6pt}AIML Engineering Lab}}}

\vspace{6pt}

{\fontsize{9}{11}\selectfont\color{codemuted}PROJECT 010}

\vspace{1pt}
{\fontsize{22}{27}\selectfont\bfseries\color{titleblue} Hyperparameter Optimization}

\vspace{2pt}
{\normalsize\color{accent} Mastering Grid, Random, and Bayesian Optimization for ML performance.}

\vspace{4pt}
{\color{tableheader}\rule{\linewidth}{0.4pt}}

\vspace{3pt}

\section{Project Overview}

I tackled two regression problems using Grid Search, Random Search, and Bayesian Optimization (Optuna). My goal was to optimize hyperparameters for improved model performance, yielding an R² of 0.9741 on the 3D Printer dataset and 0.9733 on the PLL Loop Filter dataset. These strategies showed distinct advantages depending on the search space.

\textbf{Grid Search:} Exhaustively tests every parameter combination.

\textbf{Random Search:} Samples parameter combinations randomly for exploration.

\textbf{Bayesian Optimization:} Utilizes surrogate models for efficient searching.

\textbf{Optuna:} State-of-the-art framework for hyperparameter optimization.


%% ── DATASETS ───────────────────────────────────────────────────────────────
\needspace{8\baselineskip}
\section{Datasets}

\subsection{3D Printer Extrusion Quality}

\begin{tabularx}{\linewidth}{@{}>{\bfseries}l X @{}}
  \arrayrulecolor{tableheader}
  \toprule
  \rowcolor{tableheader} Property & Value \\
  \midrule
  \rowcolor{white} Description & Dataset focuses on the quality of 3D prints based on operational parameters. \\
  \rowcolor{tablerow} Total Rows & 5,000 \\
  \rowcolor{white} Columns & 4 (3 features + 1 target) \\
  \rowcolor{tablerow} Features & Temperature, Speed, Layer Height \\
  \rowcolor{white} Class Balance & Continuous target, 0-1 scale \\
  \rowcolor{tablerow} Train / Test Split & 80/20 \\
  \bottomrule
\end{tabularx}

\vspace{6pt}

\subsection{PLL Loop Filter Tuning}

\begin{tabularx}{\linewidth}{@{}>{\bfseries}l X @{}}
  \arrayrulecolor{tableheader}
  \toprule
  \rowcolor{tableheader} Property & Value \\
  \midrule
  \rowcolor{white} Description & Dataset aimed at tuning phase-locked loop parameters post-silicon. \\
  \rowcolor{tablerow} Total Rows & 5,000 \\
  \rowcolor{white} Columns & 4 (3 features + 1 target) \\
  \rowcolor{tablerow} Features & PLL Parameters \\
  \rowcolor{white} Class Balance & Continuous, microseconds \\
  \rowcolor{tablerow} Train / Test Split & 80/20 \\
  \bottomrule
\end{tabularx}

%% ── METHODOLOGY ────────────────────────────────────────────────────────────
\section{Methodology}

\subsection{Algorithms Used}

\begin{tabularx}{\linewidth}{@{}>{\bfseries}l X @{}}
  \arrayrulecolor{tableheader}
  \toprule
  \rowcolor{tableheader} Algorithm & Description \\
  \midrule
  \rowcolor{white} Grid Search & Chosen for exhaustive parameter testing. \\
  \rowcolor{tablerow} Random Search & Ideal for quick explorations of larger spaces. \\
  \rowcolor{white} Bayesian Optimization & Utilizes past evaluations for intelligent sampling. \\
  \rowcolor{tablerow} Optuna & Used for advanced Bayesian optimization and pruning. \\
  \bottomrule
\end{tabularx}

I loaded the data, established baselines, applied each optimization technique, and evaluated model performance using R² and RMSE metrics.

%% ── RESULTS (images mixed in) ──────────────────────────────────────────────
\section{Results}

The Bayesian Optimization approach via Optuna outperformed others, significantly improving the output quality by efficiently navigating large parameter spaces. I found that random sampling provided satisfactory results but lacked the precision achieved with Bayesian methods.

\vspace{4pt}

\begin{tabularx}{\linewidth}{@{}>{\bfseries}l X @{}}
  \arrayrulecolor{tableheader}
  \toprule
  \rowcolor{tableheader} Metric & Value \\
  \midrule
  \rowcolor{white} R² & 0.9741 \\
  \rowcolor{tablerow} RMSE & 0.0295 \\
  \rowcolor{white} R² & 0.9733 \\
  \rowcolor{tablerow} RMSE & 5.0692 \\
  \bottomrule
\end{tabularx}

\vspace{4pt}
\begin{resultbox}
  \textbf{\color{accent}Best Result}\quad {\bfseries\color{metricgreen}R² = 0.9741}
\end{resultbox}

\vspace{6pt}

\begin{figure}[H]
  \centering
  \includegraphics[width=\linewidth]{/Users/rajendarmuddasani/AIML/55_LinkedIn/010_hyperparameter_optimization/assets/wow_3d_pll_design_space.png}
\end{figure}

\begin{figure}[H]
  \centering
  \includegraphics[width=\linewidth]{/Users/rajendarmuddasani/AIML/55_LinkedIn/010_hyperparameter_optimization/assets/wow_3d_hp_landscape.png}
\end{figure}

\begin{figure}[H]
  \centering
  \includegraphics[width=\linewidth]{/Users/rajendarmuddasani/AIML/55_LinkedIn/010_hyperparameter_optimization/assets/wow_model_heatmap.png}
\end{figure}

\begin{figure}[H]
  \centering
  \includegraphics[width=\linewidth]{/Users/rajendarmuddasani/AIML/55_LinkedIn/010_hyperparameter_optimization/assets/wow_convergence_comparison.png}
\end{figure}

\begin{figure}[H]
  \centering
  \includegraphics[width=\linewidth]{/Users/rajendarmuddasani/AIML/55_LinkedIn/010_hyperparameter_optimization/assets/fig3_search_landscape.png}
\end{figure}

\begin{figure}[H]
  \centering
  \includegraphics[width=\linewidth]{/Users/rajendarmuddasani/AIML/55_LinkedIn/010_hyperparameter_optimization/assets/fig1_search_comparison.png}
\end{figure}

\begin{figure}[H]
  \centering
  \includegraphics[width=\linewidth]{/Users/rajendarmuddasani/AIML/55_LinkedIn/010_hyperparameter_optimization/assets/fig5_pll_optimization.png}
\end{figure}

\begin{figure}[H]
  \centering
  \includegraphics[width=\linewidth]{/Users/rajendarmuddasani/AIML/55_LinkedIn/010_hyperparameter_optimization/assets/fig4_cross_validation.png}
\end{figure}

\begin{figure}[H]
  \centering
  \includegraphics[width=\linewidth]{/Users/rajendarmuddasani/AIML/55_LinkedIn/010_hyperparameter_optimization/assets/pll_hpo_comparison.png}
\end{figure}

\begin{figure}[H]
  \centering
  \includegraphics[width=\linewidth]{/Users/rajendarmuddasani/AIML/55_LinkedIn/010_hyperparameter_optimization/assets/hpo_comparison.png}
\end{figure}

\begin{figure}[H]
  \centering
  \includegraphics[width=\linewidth]{/Users/rajendarmuddasani/AIML/55_LinkedIn/010_hyperparameter_optimization/assets/optuna_history.png}
\end{figure}

\begin{figure}[H]
  \centering
  \includegraphics[width=\linewidth]{/Users/rajendarmuddasani/AIML/55_LinkedIn/010_hyperparameter_optimization/assets/fig2_optuna_insights.png}
\end{figure}

\begin{figure}[H]
  \centering
  \includegraphics[width=\linewidth]{/Users/rajendarmuddasani/AIML/55_LinkedIn/010_hyperparameter_optimization/assets/param_importance.png}
\end{figure}

The visualizations illustrate the convergence rates of each optimization method and parameter significance across datasets.

%% ── TECH STACK & KEY LEARNINGS ──────────────────────────────────────────────
\section{Tech Stack \& Key Learnings}

\begin{tabularx}{\linewidth}{@{}>{\bfseries}l X @{}}
  \arrayrulecolor{tableheader}
  \toprule
  \rowcolor{tableheader} Info & Detail \\
  \midrule
  \rowcolor{white} GitHub & \href{https://github.com/AIML-Engineering-Lab/010_hyperparameter_optimization}{\nolinkurl{https://github.com/AIML-Engineering-Lab/010_hyperparameter_optimization}} \\
  \rowcolor{tablerow} Tech Stack & Python 3.11, scikit-learn, pandas, Matplotlib \\
  \rowcolor{white} Licence & MIT \\
  \bottomrule
\end{tabularx}

\vspace{4pt}
\begin{itemize}[leftmargin=18pt,itemsep=2pt]
  \item Grid Search is effective for small spaces.
  \item Random Search is faster and simpler for medium spaces.
  \item Bayesian Optimization significantly optimizes performance in larger search spaces.
\end{itemize}

\end{document}
